\documentclass[12pt,a4paper]{article}
\usepackage[utf8]{inputenc}
\usepackage[T1]{fontenc}
\usepackage{lmodern}
\usepackage[margin=1in]{geometry}
\usepackage{amsmath,amssymb}
\usepackage{xcolor}
\usepackage{listings}
\usepackage{booktabs}
\usepackage{enumitem}
\usepackage{tikz}
\usetikzlibrary{shapes,arrows,positioning}
\usepackage{tcolorbox}
\tcbuselibrary{skins,breakable}
\usepackage{fancyhdr}
\usepackage{hyperref}
\usepackage{array}
\usepackage{multirow}

\definecolor{codegreen}{rgb}{0,0.6,0}
\definecolor{codegray}{rgb}{0.5,0.5,0.5}
\definecolor{codepurple}{rgb}{0.58,0,0.82}
\definecolor{backcolour}{rgb}{0.95,0.95,0.97}
\definecolor{accentblue}{RGB}{41,98,255}
\definecolor{accentred}{RGB}{220,50,47}
\definecolor{accentgreen}{RGB}{0,150,80}
\definecolor{lightblue}{RGB}{230,240,255}
\definecolor{lightyellow}{RGB}{255,250,230}
\definecolor{lightgreen}{RGB}{230,255,235}
\definecolor{lightred}{RGB}{255,235,235}

\lstdefinestyle{mystyle}{
    backgroundcolor=\color{backcolour}, commentstyle=\color{codegreen},
    keywordstyle=\color{accentblue}\bfseries, numberstyle=\tiny\color{codegray},
    stringstyle=\color{codepurple}, basicstyle=\ttfamily\small,
    breaklines=true, numbers=left, numbersep=5pt, frame=single,
    rulecolor=\color{codegray}, tabsize=2
}
\lstset{style=mystyle}

\lstdefinelanguage{SQL}{
    keywords={SELECT,FROM,WHERE,INSERT,INTO,VALUES,UPDATE,SET,DELETE,CREATE,TABLE,ALTER,DROP,JOIN,INNER,LEFT,RIGHT,FULL,OUTER,ON,AND,OR,NOT,IN,EXISTS,BETWEEN,LIKE,ORDER,BY,GROUP,HAVING,DISTINCT,AS,UNION,ALL,PRIMARY,KEY,FOREIGN,REFERENCES,INDEX,VIEW,TRIGGER,BEGIN,END,COMMIT,ROLLBACK,SAVEPOINT,GRANT,REVOKE,CASCADE,UNIQUE,CHECK,DEFAULT,AUTO_INCREMENT,LIMIT,OFFSET,ASC,DESC,WITH,EXCEPT,INTERSECT,COUNT,SUM,AVG,MAX,MIN,FOR,PROCEDURE,FUNCTION,CALL,AFTER,EACH,ROW,OLD,NEW,NOW,LAST_INSERT_ID,CURRENT_USER_ID,OVER,PARTITION,DENSE_RANK,TRUNCATE,NULL,IS},
    sensitive=false, morecomment=[l]{--}, morestring=[b]'
}

\newtcolorbox{conceptbox}[1]{colback=lightblue, colframe=accentblue, fonttitle=\bfseries\large, title=#1, breakable, sharp corners}
\newtcolorbox{examplebox}[1]{colback=lightgreen, colframe=accentgreen, fonttitle=\bfseries\large, title=#1, breakable, sharp corners}
\newtcolorbox{warningbox}[1]{colback=lightred, colframe=accentred, fonttitle=\bfseries\large, title=#1, breakable, sharp corners}
\newtcolorbox{tipbox}[1]{colback=lightyellow, colframe=orange, fonttitle=\bfseries\large, title=#1, breakable, sharp corners}

\pagestyle{fancy}
\fancyhf{}
\fancyhead[L]{\textbf{DBMS Concepts}}
\fancyhead[R]{\textbf{Interview Guide}}
\fancyfoot[C]{\thepage}
\renewcommand{\headrulewidth}{1pt}
\hypersetup{colorlinks=true, linkcolor=accentblue, urlcolor=accentblue}

\title{
    \vspace{-1cm}
    {\Huge\textbf{DBMS Concepts}}\\[0.3cm]
    {\LARGE Complete Interview Preparation Guide}\\[0.3cm]
    {\large With Real-World Examples, SQL, Diagrams \& Theory}
}
\author{}
\date{\today}

\begin{document}
\maketitle
\thispagestyle{fancy}
\tableofcontents
\newpage

\section{Introduction to DBMS}

\begin{conceptbox}{What is a DBMS?}
A \textbf{Database Management System} is software that stores, retrieves, and manages data efficiently, providing an interface between the database and users/applications.
\end{conceptbox}

\subsection{File System vs DBMS}
\begin{center}
\begin{tabular}{@{}l p{4.5cm} p{4.5cm}@{}}
\toprule
\textbf{Feature} & \textbf{File System} & \textbf{DBMS} \\
\midrule
Redundancy & High (duplicate files) & Low (normalization) \\
Consistency & Hard to maintain & Enforced by constraints \\
Concurrent Access & No support & Transactions \& locking \\
Security & OS-level only & Fine-grained (row/column) \\
Recovery & Manual & Automatic (WAL, logs) \\
Query Language & None & SQL \\
\bottomrule
\end{tabular}
\end{center}

\subsection{Types of DBMS}
\begin{center}
\begin{tabular}{@{}l l l@{}}
\toprule
\textbf{Type} & \textbf{Model} & \textbf{Example} \\
\midrule
Relational & Tables (rows/columns) & MySQL, PostgreSQL, Oracle \\
Document & JSON documents & MongoDB, CouchDB \\
Key-Value & Key $\rightarrow$ Value & Redis, DynamoDB \\
Column-family & Column families & Cassandra, HBase \\
Graph & Nodes + Edges & Neo4j \\
\bottomrule
\end{tabular}
\end{center}

\subsection{Three-Schema Architecture}
\begin{itemize}[leftmargin=*]
    \item \textbf{External Level:} User views (different views for different users)
    \item \textbf{Conceptual Level:} Logical schema (tables, relationships)
    \item \textbf{Internal Level:} Physical storage (files, indexes, blocks)
\end{itemize}

\textbf{Data Independence:}
\begin{itemize}[leftmargin=*]
    \item \textbf{Logical:} Change conceptual schema without changing external views.
    \item \textbf{Physical:} Change storage without changing conceptual schema.
\end{itemize}

\subsection{DBMS Languages}
\begin{center}
\begin{tabular}{@{}l l l@{}}
\toprule
\textbf{Language} & \textbf{Purpose} & \textbf{Commands} \\
\midrule
DDL & Define schema & CREATE, ALTER, DROP \\
DML & Query/modify data & SELECT, INSERT, UPDATE, DELETE \\
DCL & Access control & GRANT, REVOKE \\
TCL & Transactions & COMMIT, ROLLBACK, SAVEPOINT \\
\bottomrule
\end{tabular}
\end{center}

\section{ER Model (Entity-Relationship)}

\begin{conceptbox}{ER Model}
A high-level conceptual model to design databases. Describes \textbf{entities} (things), \textbf{attributes} (properties), and \textbf{relationships} (associations).
\end{conceptbox}

\subsection{Components}
\begin{center}
\begin{tabular}{@{}l l l@{}}
\toprule
\textbf{Component} & \textbf{Symbol} & \textbf{Example} \\
\midrule
Entity & Rectangle & Student, Course \\
Key Attribute & \underline{Underlined} ellipse & \underline{roll\_no} \\
Multivalued Attr & Double ellipse & phone\_numbers \\
Derived Attr & Dashed ellipse & age (from DOB) \\
Composite Attr & Sub-ellipses & name $\to$ (first, last) \\
Relationship & Diamond & ``enrolls\_in'' \\
Weak Entity & Double rectangle & Dependent (of Employee) \\
\bottomrule
\end{tabular}
\end{center}

\subsection{Cardinality}
\begin{itemize}[leftmargin=*]
    \item \textbf{1:1} --- Person : Passport
    \item \textbf{1:N} --- Department : Employees
    \item \textbf{M:N} --- Students : Courses
\end{itemize}

\begin{examplebox}{ER to Tables --- University}
\begin{lstlisting}[language=SQL]
CREATE TABLE Department (
    dept_id INT PRIMARY KEY, dept_name VARCHAR(50));

CREATE TABLE Professor (
    emp_id INT PRIMARY KEY, name VARCHAR(100),
    salary DECIMAL(10,2),
    dept_id INT REFERENCES Department(dept_id));

CREATE TABLE Student (
    roll_no INT PRIMARY KEY, first_name VARCHAR(50),
    last_name VARCHAR(50), dob DATE,
    dept_id INT REFERENCES Department(dept_id));

CREATE TABLE Course (
    course_id INT PRIMARY KEY, title VARCHAR(100),
    credits INT,
    prof_id INT REFERENCES Professor(emp_id));

-- M:N junction table
CREATE TABLE Enrollment (
    roll_no INT REFERENCES Student(roll_no),
    course_id INT REFERENCES Course(course_id),
    grade CHAR(2),
    PRIMARY KEY (roll_no, course_id));

-- Weak entity
CREATE TABLE Dependent (
    emp_id INT REFERENCES Professor(emp_id),
    dep_name VARCHAR(50), relation VARCHAR(20),
    PRIMARY KEY (emp_id, dep_name));
\end{lstlisting}
\end{examplebox}

\subsection{ER Mapping Rules}
\begin{center}
\begin{tabular}{@{}l p{8cm}@{}}
\toprule
\textbf{ER Construct} & \textbf{Rule} \\
\midrule
Strong Entity & Table with PK = key attribute \\
Weak Entity & Table with PK = partial key + owner's PK \\
1:1 Relationship & FK on total participation side \\
1:N Relationship & FK on the N-side \\
M:N Relationship & New junction table with FKs \\
Multivalued Attr & Separate table (entity PK, value) \\
\bottomrule
\end{tabular}
\end{center}

\section{Relational Model \& Relational Algebra}

\subsection{Keys}
\begin{center}
\begin{tabular}{@{}l p{5cm} p{4cm}@{}}
\toprule
\textbf{Key} & \textbf{Definition} & \textbf{Example} \\
\midrule
Super Key & Any set uniquely identifying a tuple & \{roll\_no\}, \{roll\_no, name\} \\
Candidate Key & Minimal super key & \{roll\_no\}, \{email\} \\
Primary Key & Chosen candidate key & roll\_no \\
Foreign Key & References PK of another table & Student.dept\_id \\
Composite Key & PK with multiple attributes & (roll\_no, course\_id) \\
\bottomrule
\end{tabular}
\end{center}

\subsection{Integrity Constraints}
\begin{enumerate}[leftmargin=*]
    \item \textbf{Domain:} Values from valid domain (age $> 0$)
    \item \textbf{Key:} No duplicate key values
    \item \textbf{Entity Integrity:} PK cannot be NULL
    \item \textbf{Referential Integrity:} FK must match PK or be NULL
\end{enumerate}

\subsection{Relational Algebra}
\begin{center}
\begin{tabular}{@{}l l p{5cm}@{}}
\toprule
\textbf{Operator} & \textbf{Symbol} & \textbf{Description} \\
\midrule
Selection & $\sigma_{cond}(R)$ & Filter rows \\
Projection & $\pi_{cols}(R)$ & Select columns \\
Union & $R \cup S$ & All tuples in R or S \\
Difference & $R - S$ & Tuples in R not in S \\
Cartesian Product & $R \times S$ & All combinations \\
Natural Join & $R \bowtie S$ & Join on common attrs \\
Division & $R \div S$ & Tuples in R with ALL of S \\
\bottomrule
\end{tabular}
\end{center}

\begin{examplebox}{Relational Algebra Examples}
Find CS students: $\pi_{\text{name}}(\sigma_{\text{dept}='CS'}(\text{Student}))$

Find students in ``DBMS'': $\pi_{\text{name}}(\text{Student} \bowtie \text{Enrollment} \bowtie \sigma_{\text{title}='DBMS'}(\text{Course}))$

Students enrolled in ALL courses: $\pi_{\text{roll,cid}}(\text{Enrollment}) \div \pi_{\text{cid}}(\text{Course})$
\end{examplebox}

\section{SQL --- Structured Query Language}

\subsection{DDL}
\begin{lstlisting}[language=SQL]
CREATE TABLE Employee (
    emp_id INT PRIMARY KEY AUTO_INCREMENT,
    name VARCHAR(100) NOT NULL,
    email VARCHAR(100) UNIQUE,
    salary DECIMAL(10,2) DEFAULT 50000,
    dept_id INT REFERENCES Department(dept_id)
        ON DELETE SET NULL ON UPDATE CASCADE,
    CHECK (salary > 0));

ALTER TABLE Employee ADD COLUMN phone VARCHAR(15);
DROP TABLE Employee;      -- Removes table + schema
TRUNCATE TABLE Employee;  -- Removes all rows, keeps schema
\end{lstlisting}

\begin{tipbox}{DROP vs TRUNCATE vs DELETE}
\begin{center}
\begin{tabular}{@{}l l l l@{}}
\toprule
& \textbf{DROP} & \textbf{TRUNCATE} & \textbf{DELETE} \\
\midrule
Type & DDL & DDL & DML \\
Removes & Table+data+schema & All rows & Specific rows \\
Rollback? & No & No & Yes \\
Triggers? & No & No & Yes \\
\bottomrule
\end{tabular}
\end{center}
\end{tipbox}

\subsection{DML}
\begin{lstlisting}[language=SQL]
INSERT INTO Employee (name, salary, dept_id) VALUES
    ('Harsh', 90000, 1), ('Amit', 75000, 2);

UPDATE Employee SET salary = salary * 1.10 WHERE dept_id = 1;
DELETE FROM Employee WHERE emp_id = 5;
\end{lstlisting}

\subsection{SELECT --- Real-World E-commerce Queries}

\begin{examplebox}{E-commerce Database Queries}
Tables: Customer(cust\_id, name, city), Product(prod\_id, name, price, category), Orders(order\_id, cust\_id, order\_date, total), OrderItem(order\_id, prod\_id, quantity, price)

\textbf{Q1: Mumbai customers who spent > 50K:}
\begin{lstlisting}[language=SQL]
SELECT c.name, SUM(o.total) AS total_spent
FROM Customer c JOIN Orders o ON c.cust_id = o.cust_id
WHERE c.city = 'Mumbai'
GROUP BY c.cust_id, c.name
HAVING SUM(o.total) > 50000
ORDER BY total_spent DESC;
\end{lstlisting}

\textbf{Q2: Second highest price:}
\begin{lstlisting}[language=SQL]
SELECT DISTINCT price FROM Product
ORDER BY price DESC LIMIT 1 OFFSET 1;

-- Or using DENSE_RANK:
SELECT name, price FROM (
    SELECT name, price,
        DENSE_RANK() OVER (ORDER BY price DESC) AS rnk
    FROM Product) t WHERE rnk = 2;
\end{lstlisting}

\textbf{Q3: Customers who bought every Electronics product (Division):}
\begin{lstlisting}[language=SQL]
SELECT c.name FROM Customer c
WHERE NOT EXISTS (
    SELECT prod_id FROM Product WHERE category='Electronics'
    EXCEPT
    SELECT oi.prod_id FROM Orders o
    JOIN OrderItem oi ON o.order_id = oi.order_id
    WHERE o.cust_id = c.cust_id);
\end{lstlisting}

\textbf{Q4: Running total per customer:}
\begin{lstlisting}[language=SQL]
SELECT cust_id, order_date, total,
    SUM(total) OVER (PARTITION BY cust_id
        ORDER BY order_date) AS running_total
FROM Orders;
\end{lstlisting}
\end{examplebox}

\subsection{Joins}
\begin{center}
\begin{tabular}{@{}l p{8cm}@{}}
\toprule
\textbf{Type} & \textbf{Result} \\
\midrule
INNER JOIN & Only matching rows \\
LEFT JOIN & All left + matching right (NULL if no match) \\
RIGHT JOIN & All right + matching left \\
FULL OUTER JOIN & All from both (NULL where no match) \\
CROSS JOIN & Cartesian product \\
SELF JOIN & Table joined with itself \\
\bottomrule
\end{tabular}
\end{center}

\begin{examplebox}{Join with Data}
Employee: (1,Harsh,10), (2,Amit,20), (3,Priya,NULL) \\
Department: (10,Engineering), (20,Marketing), (30,Finance)

\textbf{INNER:} (Harsh,Engg), (Amit,Mktg) --- Priya and Finance excluded\\
\textbf{LEFT:} (Harsh,Engg), (Amit,Mktg), (Priya,NULL) --- all employees\\
\textbf{FULL OUTER:} adds (NULL,Finance) too
\end{examplebox}

\subsection{Subqueries}
\begin{examplebox}{Correlated Subquery: Above dept average salary}
\begin{lstlisting}[language=SQL]
SELECT name, salary, dept_id FROM Employee e1
WHERE salary > (
    SELECT AVG(salary) FROM Employee e2
    WHERE e2.dept_id = e1.dept_id);
-- Executes inner query ONCE PER ROW of outer query
\end{lstlisting}
\end{examplebox}

\subsection{Views, Triggers, Procedures}
\begin{lstlisting}[language=SQL]
CREATE VIEW HighEarners AS
SELECT name, salary FROM Employee WHERE salary > 100000;

CREATE TRIGGER salary_audit AFTER UPDATE ON Employee
FOR EACH ROW BEGIN
    INSERT INTO AuditLog VALUES(OLD.emp_id, OLD.salary, NEW.salary, NOW());
END;

CREATE PROCEDURE GiveRaise(IN dept INT, IN pct DECIMAL)
BEGIN
    UPDATE Employee SET salary=salary*(1+pct/100) WHERE dept_id=dept;
END;
CALL GiveRaise(10, 15);
\end{lstlisting}

\section{Normalization}

\begin{conceptbox}{Normalization}
Process of organizing a database to reduce \textbf{redundancy} and prevent \textbf{anomalies} (insertion, deletion, update). Decomposes tables following \textbf{normal forms}.
\end{conceptbox}

\subsection{Anomalies}
\begin{warningbox}{Why Normalize?}
Table: StudentCourse(roll, name, dept, course, prof)

\begin{itemize}[leftmargin=*]
    \item \textbf{Update:} Change Harsh's dept $\to$ update multiple rows
    \item \textbf{Insertion:} Can't add dept without a student
    \item \textbf{Deletion:} Drop last student in a course $\to$ lose prof info
\end{itemize}
\end{warningbox}

\subsection{Functional Dependencies}
$X \to Y$: same X value $\Rightarrow$ same Y value. E.g., roll $\to$ name.

\textbf{Armstrong's Axioms:}
Reflexivity ($Y \subseteq X \Rightarrow X \to Y$), Augmentation ($X \to Y \Rightarrow XZ \to YZ$), Transitivity ($X \to Y, Y \to Z \Rightarrow X \to Z$).

\textbf{Derived:} Union ($X \to Y, X \to Z \Rightarrow X \to YZ$), Decomposition ($X \to YZ \Rightarrow X \to Y$).

\begin{examplebox}{Finding Candidate Keys via Closure}
$R(A,B,C,D,E)$, FDs: $A \to B$, $BC \to D$, $D \to E$, $E \to A$

$AC^+ = \{A,C\} \xrightarrow{A \to B} \{A,B,C\} \xrightarrow{BC \to D} \{A,B,C,D\} \xrightarrow{D \to E} \{A,B,C,D,E\} = R$ $\checkmark$

$AC$ is a candidate key. $C$ never on RHS $\to$ must be in every CK.

All candidate keys: $\{AC, BC, DC, EC\}$.
\end{examplebox}

\subsection{Normal Forms}

\subsubsection{1NF --- Atomic values only}
No multivalued attributes. E.g., phones = ``9876, 9123'' $\to$ separate rows or table.

\subsubsection{2NF --- No partial dependencies}
Non-prime attribute must not depend on subset of candidate key.

\begin{examplebox}{2NF Fix}
\texttt{Enrollment(\underline{roll, course\_id}, student\_name, grade)}\\
roll $\to$ student\_name (partial dep on PK) $\Rightarrow$ Decompose:\\
Student(\underline{roll}, student\_name) + Enrollment(\underline{roll, course\_id}, grade)
\end{examplebox}

\subsubsection{3NF --- No transitive dependencies}
For every FD $X \to A$: either $X$ is superkey OR $A$ is prime attribute.

\begin{examplebox}{3NF Fix}
\texttt{Employee(\underline{emp\_id}, dept\_id, dept\_name)}\\
emp\_id $\to$ dept\_id $\to$ dept\_name (transitive) $\Rightarrow$ Decompose:\\
Employee(\underline{emp\_id}, dept\_id) + Department(\underline{dept\_id}, dept\_name)
\end{examplebox}

\subsubsection{BCNF --- Every determinant is a superkey}
Stricter than 3NF. For every FD $X \to Y$, $X$ must be a superkey (no exception for prime attrs).

\begin{examplebox}{BCNF Violation}
\texttt{Teaching(\underline{student, course}, professor)}\\
FDs: \{student,course\} $\to$ professor, professor $\to$ course\\
professor is not a superkey $\Rightarrow$ violates BCNF.\\
Fix: ProfCourse(\underline{professor}, course) + StudentProf(\underline{student, professor})
\end{examplebox}

\subsection{Normal Forms Summary}
\begin{center}
\begin{tabular}{@{}l p{4.5cm} p{4cm}@{}}
\toprule
\textbf{NF} & \textbf{Rule} & \textbf{Eliminates} \\
\midrule
1NF & Atomic values & Repeating groups \\
2NF & No partial deps & Partial key redundancy \\
3NF & No transitive deps & Non-key redundancy \\
BCNF & Determinant = superkey & All FD redundancy \\
4NF & No multi-valued deps & Independent MVDs \\
\bottomrule
\end{tabular}
\end{center}

$\text{1NF} \subset \text{2NF} \subset \text{3NF} \subset \text{BCNF} \subset \text{4NF}$. Most systems aim for \textbf{3NF or BCNF}.

\subsection{Lossless Join \& Dependency Preservation}
Decomposition of $R$ into $R_1, R_2$ is \textbf{lossless} iff:
$R_1 \cap R_2 \to R_1$ or $R_1 \cap R_2 \to R_2$ (common attrs form superkey of one table).

\textbf{3NF decomposition:} always lossless + dependency preserving.\\
\textbf{BCNF decomposition:} always lossless, but may NOT preserve dependencies.

\section{Transactions \& ACID Properties}

\begin{conceptbox}{Transaction}
A logical unit of work: one or more SQL operations that either \textbf{all succeed} or \textbf{all fail}.
\end{conceptbox}

\subsection{ACID}
\begin{center}
\begin{tabular}{@{}l p{4cm} p{5cm}@{}}
\toprule
\textbf{Property} & \textbf{Meaning} & \textbf{Real-World Example} \\
\midrule
Atomicity & All or nothing & Bank transfer: debit+credit both or neither \\
Consistency & Valid state $\to$ valid state & Total money unchanged after transfer \\
Isolation & No interference & Two bookings for last seat: one wins \\
Durability & Committed = permanent & Power failure after commit $\to$ data safe \\
\bottomrule
\end{tabular}
\end{center}

\begin{examplebox}{Bank Transfer Transaction}
\begin{lstlisting}[language=SQL]
BEGIN TRANSACTION;
UPDATE Account SET balance = balance - 10000 WHERE acc_id = 'A';
-- If crash here: Atomicity ensures rollback
UPDATE Account SET balance = balance + 10000 WHERE acc_id = 'B';
COMMIT;  -- Durability: permanent even if crash after this
\end{lstlisting}
\end{examplebox}

\subsection{Schedules \& Serializability}

Two operations \textbf{conflict} if: different transactions + same data item + at least one is write.

\textbf{Conflict serializable:} Build precedence graph. \textbf{No cycle} $\Rightarrow$ conflict serializable.

\begin{examplebox}{Precedence Graph Test}
T1: R(A), W(A), R(B), W(B) \quad T2: R(A), W(A), R(B), W(B)

If T1.R(A) before T2.W(A): edge T1$\to$T2\\
If T2.R(A) before T1.W(A): edge T2$\to$T1\\
Both edges $\Rightarrow$ \textbf{Cycle} $\Rightarrow$ NOT conflict serializable.
\end{examplebox}

\subsection{Recoverability}
$\text{Serial} \subset \text{Strict} \subset \text{Cascadeless} \subset \text{Recoverable}$

\textbf{Cascadeless:} Only read committed data. \textbf{Strict:} Don't touch item until writer commits/aborts.

\section{Concurrency Control}

\subsection{Problems Without Concurrency Control}
\begin{itemize}[leftmargin=*]
    \item \textbf{Dirty Read:} Read uncommitted data
    \item \textbf{Non-repeatable Read:} Same query, different results
    \item \textbf{Phantom Read:} New rows appear
    \item \textbf{Lost Update:} Two writes, one lost
\end{itemize}

\subsection{Isolation Levels}
\begin{center}
\begin{tabular}{@{}l c c c@{}}
\toprule
\textbf{Level} & \textbf{Dirty} & \textbf{Non-repeat} & \textbf{Phantom} \\
\midrule
Read Uncommitted & Yes & Yes & Yes \\
Read Committed & No & Yes & Yes \\
Repeatable Read & No & No & Yes \\
Serializable & No & No & No \\
\bottomrule
\end{tabular}
\end{center}

\subsection{Two-Phase Locking (2PL)}
\begin{conceptbox}{2PL}
\textbf{Growing phase:} acquire locks, never release. \textbf{Shrinking phase:} release locks, never acquire. Guarantees \textbf{conflict serializability}. Can cause \textbf{deadlocks}.

Variants: \textbf{Basic} (not recoverable), \textbf{Strict} (hold X-locks till commit), \textbf{Rigorous} (hold ALL locks till commit).
\end{conceptbox}

\textbf{Lock compatibility:} S-S $\checkmark$, S-X $\times$, X-X $\times$.

\subsection{Deadlock Handling}
\begin{itemize}[leftmargin=*]
    \item \textbf{Wait-Die:} Older waits, younger dies (rollback)
    \item \textbf{Wound-Wait:} Older wounds (forces rollback of younger), younger waits
    \item \textbf{Timeout:} Rollback if waiting too long
    \item \textbf{Detection:} Wait-for graph; cycle $\to$ rollback victim
\end{itemize}

\subsection{Timestamp-Based Protocol}
Each transaction gets timestamp at start. Rules ensure equivalent serial order by timestamp. \textbf{Deadlock-free} but may cause cascading rollbacks.

\subsection{MVCC (Multi-Version Concurrency Control)}
Maintain multiple versions. Readers see snapshot at their start time. \textbf{Readers never block writers.} Used by PostgreSQL, MySQL InnoDB, Oracle.

\section{Indexing}

\begin{conceptbox}{Why Indexing?}
Without index: full scan $O(n)$. With index: $O(\log n)$ or $O(1)$.
\end{conceptbox}

\subsection{Index Types}
\begin{center}
\begin{tabular}{@{}l l l@{}}
\toprule
\textbf{Type} & \textbf{On} & \textbf{Example} \\
\midrule
Primary & Ordering key of sorted file & Index on emp\_id \\
Clustering & Non-key ordering field & Index on dept\_id \\
Secondary & Non-ordering field & Index on name \\
Dense & Entry for every record & Fast, more space \\
Sparse & Entry per block & Less space, slower \\
\bottomrule
\end{tabular}
\end{center}

\subsection{B+ Tree (Most Important)}
\begin{conceptbox}{B+ Tree Properties}
\begin{itemize}
    \item Data stored \textbf{only in leaf nodes}. Internal nodes = routing keys.
    \item \textbf{Leaf nodes linked} $\to$ efficient range queries.
    \item Balanced: all leaves at same depth.
    \item Node with $k$ children has $k-1$ keys.
    \item Internal: min $\lceil m/2 \rceil$ children. Max $m$ children.
\end{itemize}
\end{conceptbox}

\begin{examplebox}{B+ Tree Calculation}
Block = 4096B, Key = 8B, Pointer = 8B.

Internal order: $16p - 8 \leq 4096 \Rightarrow p = 256$ (255 keys/node).

For 1,000,000 records: Height $\approx 3$. Only \textbf{3 disk reads} to find any record!
\end{examplebox}

\subsection{Hashing vs B+ Tree}
\begin{center}
\begin{tabular}{@{}l l l@{}}
\toprule
& \textbf{Hashing} & \textbf{B+ Tree} \\
\midrule
Exact match & $O(1)$ & $O(\log n)$ \\
Range query & Not supported & Excellent \\
Ordered access & No & Yes \\
\bottomrule
\end{tabular}
\end{center}

\section{Query Processing \& Optimization}

\textbf{Steps:} SQL $\to$ Parser $\to$ Relational Algebra $\to$ Optimizer $\to$ Execution $\to$ Result.

\subsection{Join Algorithms}
\begin{center}
\begin{tabular}{@{}l l l@{}}
\toprule
\textbf{Algorithm} & \textbf{Cost} & \textbf{When} \\
\midrule
Nested Loop & $O(nm)$ & Small tables \\
Sort-Merge & $O(n\log n + m\log m)$ & Sortable data \\
Hash Join & $O(n+m)$ & Equality joins \\
Index Nested Loop & $O(n\log m)$ & Index on join col \\
\bottomrule
\end{tabular}
\end{center}

\textbf{Heuristics:} Push selections down, push projections down, use indexes, avoid Cartesian products.

\section{Recovery System}

\subsection{Write-Ahead Log (WAL)}
Log every operation \textbf{before} modifying DB. Format: $\langle T_i, X, \text{old}, \text{new} \rangle$.

\textbf{Undo:} Transaction didn't commit $\to$ restore old values.\\
\textbf{Redo:} Transaction committed but not flushed $\to$ apply new values.

\subsection{Checkpointing}
Marks point where committed changes are on disk. Recovery only processes after last checkpoint.

\begin{examplebox}{Recovery Example}
Log: T1 start $\to$ T1 writes A $\to$ T2 start $\to$ CHECKPOINT $\to$ T2 writes B $\to$ T1 commit $\to$ T3 start $\to$ T3 writes C $\to$ CRASH

Recovery: T1 committed $\Rightarrow$ \textbf{REDO}. T2 not committed $\Rightarrow$ \textbf{UNDO}. T3 not committed $\Rightarrow$ \textbf{UNDO}.
\end{examplebox}

\section{Real-World Case Study: Amazon E-Commerce}

\begin{conceptbox}{How DBMS Powers Amazon}
Traces a customer purchase showing every DBMS concept in practice.
\end{conceptbox}

\subsection{Database Design}
Entities: Customer, Product, Order, OrderItem, Seller, Review, Payment.

Relationships: Customer places Order (1:N), Order contains Products (M:N via OrderItem), Seller lists Products (1:N), Customer writes Reviews (M:N).

After \textbf{3NF normalization}: 7 tables, each fact stored once.

\subsection{Transaction: Payment Processing}
\begin{lstlisting}[language=SQL]
BEGIN TRANSACTION;
SELECT quantity FROM Inventory WHERE prod_id=101 FOR UPDATE;
UPDATE Inventory SET quantity = quantity-1 WHERE prod_id=101;
INSERT INTO Orders (cust_id,order_date,total,status)
    VALUES (5001, NOW(), 79999, 'PENDING');
INSERT INTO Payment (order_id,method,amount,status)
    VALUES (LAST_INSERT_ID(), 'UPI', 79999, 'SUCCESS');
UPDATE Orders SET status='CONFIRMED'
    WHERE order_id=LAST_INSERT_ID();
COMMIT;
-- If payment fails: ROLLBACK (atomicity restores stock)
\end{lstlisting}

\subsection{Concurrency: Flash Sale}
10,000 users, last 5 iPhones:
\begin{lstlisting}[language=SQL]
-- Without control: LOST UPDATE (both read qty=5, both set 4)
-- With 2PL: SELECT ... FOR UPDATE locks the row
-- User B waits until User A commits
-- Amazon also uses: Redis atomic DECR, queue-based ordering
\end{lstlisting}

\subsection{Indexing: Why Search is Fast}
\begin{lstlisting}[language=SQL]
CREATE INDEX idx_prod_title ON Product(title);
CREATE INDEX idx_prod_cat_price ON Product(category, price);
CREATE INDEX idx_order_cust ON Orders(cust_id);
-- Without index: scan 10M rows. With B+ Tree: 3-4 disk reads!
\end{lstlisting}

\subsection{Recovery: Crash During Order}
WAL: T start $\to$ Inventory update $\to$ Order insert $\to$ CRASH.

T never committed $\Rightarrow$ UNDO: restore inventory, delete order. Customer: ``Order failed, retry.''

\section{Important Interview Questions}

\begin{enumerate}[leftmargin=*, label=\textbf{Q\arabic*.}]

\item \textbf{ACID properties?} Atomicity (all-or-nothing), Consistency (valid states), Isolation (no interference), Durability (permanent after commit).

\item \textbf{Difference between DELETE, TRUNCATE, DROP?} DELETE: DML, row-by-row, rollback-able. TRUNCATE: DDL, fast, no rollback. DROP: removes entire table.

\item \textbf{What is normalization?} Decomposing tables to reduce redundancy. 1NF (atomic) $\to$ 2NF (no partial dep) $\to$ 3NF (no transitive dep) $\to$ BCNF (determinant=superkey).

\item \textbf{Difference between 3NF and BCNF?} 3NF allows FD $X \to A$ if $A$ is prime. BCNF requires $X$ to be superkey always. BCNF may not preserve dependencies.

\item \textbf{What is a B+ Tree?} Balanced tree with data only in leaves, leaves linked. Supports range queries. Height $\sim 3$ for millions of records.

\item \textbf{What is 2PL?} Two-Phase Locking: growing (acquire) then shrinking (release). Ensures conflict serializability. Can cause deadlocks.

\item \textbf{What is MVCC?} Multi-Version Concurrency Control. Multiple versions of data. Readers see snapshots. Readers don't block writers. Used by PostgreSQL, MySQL.

\item \textbf{What is a deadlock in DBMS?} T1 holds A, wants B. T2 holds B, wants A. Solutions: Wait-Die, Wound-Wait, timeout, detection.

\item \textbf{Types of joins?} INNER (matching only), LEFT (all left), RIGHT (all right), FULL OUTER (all), CROSS (cartesian), SELF (with itself).

\item \textbf{What is a view?} Virtual table defined by a query. Doesn't store data. Provides security (show limited columns), simplicity (complex query as table).

\item \textbf{Primary key vs Unique key?} PK: no NULL, only one per table. Unique: allows NULL, multiple per table. Both enforce uniqueness.

\item \textbf{Clustered vs Non-clustered index?} Clustered: reorders table data (one per table). Non-clustered: separate structure pointing to data (multiple per table).

\item \textbf{What is WAL?} Write-Ahead Logging. Log changes before writing to DB. Enables undo (rollback) and redo (replay committed changes after crash).

\item \textbf{Conflict serializability test?} Build precedence graph from conflicting operations. No cycle = serializable. Cycle = not serializable.

\item \textbf{What is denormalization?} Adding controlled redundancy for read performance. Used in data warehouses, caching. Trade-off: faster reads, slower writes.

\item \textbf{What is a correlated subquery?} Inner query references outer query. Executes once per outer row. Slower than regular subquery.

\item \textbf{What is a foreign key?} Column referencing PK of another table. Enforces referential integrity. ON DELETE: CASCADE, SET NULL, RESTRICT.

\item \textbf{What is sharding?} Horizontal partitioning across multiple servers. Each shard holds a subset of rows. Used for scalability.

\item \textbf{SQL vs NoSQL?} SQL: structured, ACID, joins. NoSQL: flexible schema, BASE, horizontal scaling. Use SQL for transactions, NoSQL for big data/real-time.

\item \textbf{What is a trigger?} Auto-executed SQL on INSERT/UPDATE/DELETE. Used for auditing, validation, cascading updates.
\end{enumerate}

\section*{Quick Reference}

\begin{tcolorbox}[colback=lightyellow, colframe=orange, title=\textbf{Key Formulas \& Facts}]
\begin{align*}
\text{B+ Tree height} &\approx \lceil \log_p(n) \rceil \quad (p = \text{order}, n = \text{records}) \\
\text{Internal order } p &: p \times P + (p-1) \times K \leq \text{Block size} \\
\text{Leaf order } p_L &: p_L \times (K + \text{RecPtr}) + P \leq \text{Block size} \\[6pt]
\text{Lossless test} &: R_1 \cap R_2 \to R_1 \text{ or } R_1 \cap R_2 \to R_2 \\[6pt]
\text{Closure} &: X^+ = \text{all attributes determined by } X \\
\text{Candidate Key} &: X^+ = R \text{ and no proper subset of } X \text{ has this property}
\end{align*}

\textbf{Normal Form Hierarchy:} 1NF $\subset$ 2NF $\subset$ 3NF $\subset$ BCNF $\subset$ 4NF $\subset$ 5NF

\textbf{Schedule Hierarchy:} Serial $\subset$ Strict $\subset$ Cascadeless $\subset$ Recoverable $\subset$ All

\textbf{Serializability:} Conflict Ser. $\subset$ View Ser.
\end{tcolorbox}

\vfill
\begin{center}
\textit{Comprehensive DBMS guide for interview preparation. Good luck!}
\end{center}

\end{document}
